\section{Introduction}

Suspended-payload UAV transport enables aerial systems to move objects that
cannot be rigidly attached to the vehicle, but the suspended load also makes
the task sensitive to coupled vehicle-load dynamics
\cite{Sreenath2013DifferentiallyFlatHybrid,Li2023AutoTrans}. Under strong wind,
small tracking errors, payload swing, and transient speed changes can interact
with the cable-suspended payload and produce qualitatively different outcomes
across repeated runs \cite{Son2020ObstacleAvoidanceSuspendedLoad,GuerreroSanchez2017SwingAttenuation}.
Robust transport in this setting therefore requires not only feasible motion
generation and tracking, but also careful control of how aggressively the
planned motion is executed.

Payload-aware planning, MPC, and low-level control provide the core stack for
generating and tracking feasible trajectories
\cite{Li2023AutoTrans,Son2020ObstacleAvoidanceSuspendedLoad,UrbinaBrito2021PredictivePayloadTransport,Lee2010GeometricTrackingSE3}.
However, improving these components alone does not fully address the runtime
execution-aggressiveness problem. A fixed execution setting can be effective
for one mission protocol or wind/task phase while becoming too aggressive, too
conservative, or poorly timed in another. This creates a practical gap between
trajectory feasibility and balanced protocol-level robustness under repeated
strong-wind evaluation.

We study execution governance as a drop-in layer on top of an AutoTrans-like
suspended-payload UAV stack. The governor modifies execution through
\texttt{speed\_scale} and \texttt{acceleration\_scale}, so it can be attached
externally without rewriting the planner, payload MPC, or SO3 controller.
External reference or command modulation has a long control precedent
\cite{Garone2017ReferenceCommandGovernorsSurvey,Garone2016ExplicitReferenceGovernor,Li2021ActionGovernor},
but our governor is empirical and is evaluated by repeated-run behavior rather
than by formal safety certification. In the learned variant, risk-conditioned
signals are used to choose conservative or less conservative execution scales,
turning runtime aggressiveness into an interface-level adaptation problem
rather than a planner/controller redesign.

The evaluation also requires a protocol split. A single-goal mission protocol
(\texttt{goal\_repeat=1}) tests one commanded mission and is closest to nominal
single-goal execution. A goal-reissue stress protocol
(\texttt{goal\_repeat=10}) repeatedly republishes the goal and stresses
reference-update, post-arrival, and repeated-command behavior. Stress testing
and benchmark design literature cautions that protocol choice affects what a
score means \cite{Koren2018AdaptiveStressTesting,Nalic2020StressTestingScenarioBasedADS,Muller2019MeasureTargetConfusion}.
Because these protocols expose different failure modes and method rankings,
collapsing them into a mixed-protocol aggregate would obscure the trade-offs
that matter for evaluated robustness.

Our completed strong-wind evaluation shows that no current learned variant
dominates both protocols. \texttt{windlevel\_s085} is the single-goal
specialist, achieving 26/30 strict-valid runs under \texttt{goal\_repeat=1},
while \texttt{fixed\_s080} is the goal-reissue stress specialist, achieving
24/30 strict-valid runs under \texttt{goal\_repeat=10}. The learned /
risk-conditioned governor \texttt{risk\_adapter\_v1} is the most balanced
current method, with the best mean valid count (24.0/30), best worst-protocol
valid count (23/30), and lowest total regret (2). By contrast,
\texttt{risk\_adapter\_v21} remains a strong nominal / single-goal variant, but
it is weaker under stress and should not be framed as the overall best method.

This paper makes four contributions:

\begin{itemize}
  \item We introduce a drop-in learned / risk-conditioned execution governor
  for suspended-payload UAV transport that adapts \texttt{speed\_scale} and
  \texttt{acceleration\_scale} without replacing the planner, payload MPC, or
  SO3 controller.
  \item We define a protocol-split strong-wind evaluation that separates the
  single-goal mission protocol (\texttt{goal\_repeat=1}) from the goal-reissue
  stress protocol (\texttt{goal\_repeat=10}), avoiding a mixed-protocol
  aggregate as the main result.
  \item We report balanced robustness and protocol regret metrics showing that
  \texttt{risk\_adapter\_v1} is the most balanced current learned /
  risk-conditioned method, while \texttt{windlevel\_s085} and
  \texttt{fixed\_s080} are strong protocol specialists.
  \item We provide a failure-mode-aware diagnostic workflow based on
  invalid-only failure groups and representative traces, clarifying why
  aggregate strict-valid counts alone are insufficient for interpreting
  robustness.
\end{itemize}
